% ==============================================================================
% Κεφάλαιο 2: Θεωρητικό Υπόβαθρο και Σχετικές Εργασίες
% ==============================================================================

\chapter{Θεωρητικό Υπόβαθρο και Σχετικές Εργασίες}
\label{ch:background}

Η κατανόηση της αρχιτεκτονικής του προτεινόμενου συστήματος προϋποθέτει τη σύζευξη διαφορετικών πεδίων της επιστήμης υπολογιστών και της μηχανικής. Στο παρόν κεφάλαιο εξετάζονται οι θεμελιώδεις αρχές των εξελικτικών αλγορίθμων και του μοντέλου νήσων, οι κατανεμημένες δομές DHT με ευρετηρίαση RMI, οι τεχνικές προσεγγιστικής μάθησης (surrogate modeling), καθώς και οι μαθηματικές ιδιότητες των καμπυλών πλήρωσης χώρου Hilbert.
